\documentclass[11pt]{article}

\usepackage[margin=0.8in]{geometry}
\usepackage[T1]{fontenc}
\usepackage{lmodern}
\usepackage{microtype}
\usepackage{graphicx}
\usepackage{booktabs}
\usepackage{array}
\usepackage{amsmath,amssymb,tabularx}
\usepackage{xcolor}
\usepackage[font=small,labelfont=bf]{caption}
\usepackage{placeins}
\usepackage{xurl}  % long commit URLs break across lines
\usepackage[hidelinks]{hyperref}

\pagestyle{plain}
\setlength{\parindent}{0pt}
\setlength{\parskip}{5pt}
\graphicspath{{figures/}}
% Pages holding only floats fill from the top instead of spreading the floats over the page.
\makeatletter
\setlength{\@fptop}{0pt}
\setlength{\@fpsep}{14pt plus 2pt}
\setlength{\@fpbot}{0pt plus 1fil}
\makeatother
% Let figures and tables share pages with text instead of collecting on half-empty float pages.
\renewcommand{\topfraction}{0.85}
\renewcommand{\bottomfraction}{0.6}
\renewcommand{\textfraction}{0.12}
\renewcommand{\floatpagefraction}{0.8}

% Include a figure if it exists yet, otherwise a placeholder box naming it.
\newcommand{\figureorpending}[2]{\IfFileExists{figures/#1}{\includegraphics[#2]{#1}}{%
    \fbox{\parbox{0.9\linewidth}{\centering\vspace{1.5em}\textcolor{gray}{\textit{pending: \texttt{\detokenize{#1}}}}\vspace{1.5em}}}}}

\begin{document}

\begin{center}
    {\large EECS 221: Machine Learning for Science\par}
    \vspace{5pt}
    {\LARGE\bfseries
    Data--Compute Trade-offs\\[3pt]
    in Neural PDE Forecasting\par}
    \vspace{7pt}
    {\Large Suha Kagan Kose\par}
    \vspace{6pt}
    {\large Mid-Project Report\par}
\end{center}

\textbf{Project goal.} Neural PDE forecasting models must be useful under limited data and compute. We compare how
accurately and efficiently three architecture families with different inductive biases (a convolutional U-Net, a
Fourier neural operator and a vision transformer) learn Rayleigh--B\'enard convection when training data and GPU
time are limited. This check-in covers the most constrained setting of the study: 10\% of the training data and 1
or 4 hours of training on one A100 GPU. The 100\% runs follow (Section~\ref{sec:plan}).

\textbf{Milestones.} The abstract set three milestones for this check-in; all three are met.
\begin{enumerate}
    \item \textbf{Data pipeline validated against The Well's released checkpoints.} Our data path and metric
    reproduce the published test error of the released CNextU-Net and TFNO to within 1\% (Section~\ref{sec:m1}).
    \item \textbf{Budgets fixed from pilot timings.} Batch 64, fp32 and a peak learning rate of $10^{-3}$ for every
    model, from an equal-budget search; budgets count active GPU training time (Section~\ref{sec:m2}).
    \item \textbf{First 10\%-data results and a revised plan.} Six runs (3 models $\times$ 2 seeds) evaluated on
    the full test set at 1 and 4 hours, with their cost to train and to use (Sections~\ref{sec:results}--\ref{sec:plan}).
\end{enumerate}

\section{Setup}
\label{sec:setup}

\textbf{Data and task.} \texttt{rayleigh\_benard} from The Well~\cite{ohana2024well}: 1,750 simulated trajectories
of 200 frames on a $512\times128$ grid (periodic in $x$, walls in $y$), across 35 (Rayleigh, Prandtl) regimes, with
the official 80/10/10 splits. From four frames of buoyancy $b$, pressure $p$ and velocity $(u,w)$, a model predicts
the next frame, 0.25 simulation time units later~\cite{wellrb}, and forecasts further by feeding its predictions
back in (autoregressive rollout). This report trains on a nested 10\% subset: 140 trajectories, four per regime.
Windows from one trajectory are strongly correlated, so this is closer to 140 independent samples than to its
27,440 training windows.

\textbf{Models.} The Well's implementations, trained from random initialisation with matched parameter counts and
GELU activations (Table~\ref{tab:models}). They differ in how they move information across the domain, which is
what this study tests.

\begin{table}[!htbp]
\centering\footnotesize
\begin{tabularx}{\linewidth}{@{}l r >{\raggedright\arraybackslash}X >{\raggedright\arraybackslash}X@{}}
\toprule
Model & Params & How it mixes information & Inductive bias \\
\midrule
CNextU-Net~\cite{liu2022convnext} & 18.6M & ConvNeXt convolution blocks at four resolutions, joined by skip connections
& local and multi-scale \\
TFNO~\cite{li2021fno,kossaifi2024tfno} & 19.3M & global convolution through the $16\times16$ lowest Fourier modes
(factorised weights), plus pointwise layers & global and smooth: finer scales are not mixed in space \\
AViT~\cite{mccabe2024mpp} & 18.6M & $16\times16$-pixel patches as tokens, attention along each axis & global
between patches; resolution set by the patch \\
\bottomrule
\end{tabularx}
\caption{The three families. AViT's width was raised from 384 to 432 to match parameter counts.}
\label{tab:models}
\end{table}

\textbf{Training.} AdamW, one-step MSE on normalised fields, batch 64, fp32, peak learning rate $10^{-3}$ with 500
warmup steps and cosine decay to zero over four hours of active GPU training. The one-hour results are the same
runs' checkpoints after one hour (learning rate still at 85\% of its peak). Two seeds per model change the
initialisation and data order; all runs share one 10\% subset. The paper's recipe differs in learning rate,
schedule and hardware (Appendix~\ref{app:recipe}).

\textbf{Evaluation.} The error measure is The Well's VRMSE: per field $f$ and predicted frame,
\[
\mathrm{VRMSE}_f=\sqrt{\frac{\frac{1}{N}\sum_{i}(\hat y_{f,i}-y_{f,i})^2}{s_f^2+\varepsilon}},\qquad
s_f^2=\text{spatial variance of the target field},\quad \varepsilon=10^{-7},
\]
averaged over the four fields. A value of 1 means an error as large as the field's own variability. We score the
fields as stored (\emph{physical-field VRMSE}) and, to compare with the paper, also on normalised fields
(\emph{paper protocol}); the two differ only through $\varepsilon$ (Appendix~\ref{app:metrics}). One-step error
covers all 196 windows of the 175 test trajectories, split into \emph{onset} (first input frame $<50$) and
\emph{developed} convection. Rollouts start at frames 20, 80 and 140 and run 30 steps (7.5 time units); we report
the error and the fraction of rollouts that stay below VRMSE 1, 2 and 5. Horizontal kinetic-energy spectra show
which scales are right. The reference is \emph{persistence}, which copies the last input frame. Inference time is
measured at batch 1 (one forecast) and per sample at batch 64 (many forecasts at once).

\section{Milestone 1: Pipeline Validated Against the Released Checkpoints}
\label{sec:m1}

We scored The Well's released CNextU-Net and TFNO with our own data path and metric on every test window
(Table~\ref{tab:released}). Both reproduce the paper's one-step test values to within 1\%, and their rollouts land
in the paper's reported range. The released TFNO had to be run with the forward pass it was trained with: the
paper-era code skipped its Fourier layers, an error fixed upstream in 2025 (Appendix~\ref{app:tfno}).

\begin{table}[!htbp]
\centering\small
\begin{tabular}{@{}lcccc@{}}
\toprule
& \multicolumn{2}{c}{Paper-protocol one-step VRMSE} & \multicolumn{2}{c}{Rollout from frame 0, steps 6--12 / 13--30} \\
\cmidrule(lr){2-3}\cmidrule(l){4-5}
Model & Ours (released weights) & Published & Ours & Published \\
\midrule
CNextU-Net & 0.669 & 0.670 & 36 / 35 & $>10$ / $>10$ \\
TFNO & 0.663 & 0.657 & 142 / 169 & $>10$ / $>10$ \\
\bottomrule
\end{tabular}
\caption{Released checkpoints scored by our pipeline against the paper's Tables 2 and 3 (test split).}
\label{tab:released}
\end{table}

The check also showed what the headline number hides (Fig.~\ref{fig:released}). During the first $\sim$50 frames
the fluid is nearly at rest: the velocity's spatial variance, the VRMSE denominator, is close to zero, so small
absolute errors become relative errors of 1--10, while persistence is almost exact. From frame 50 on, the released
CNextU-Net scores 0.038, against 0.669 averaged over all windows. We therefore report onset and developed
convection separately.

\begin{figure}[!htbp]
\centering
\figureorpending{fig_released_by_window.pdf}{scale=1}
\caption{Released checkpoints' one-step VRMSE by window on the test set (mean over 175 trajectories), on normalised
(left) and physical fields (right). Shaded: \emph{onset} windows, first input frame below 50, a cut chosen where
the error levels off on the validation split. Persistence copies the last input frame.}
\label{fig:released}
\end{figure}

\section{Milestone 2: Budgets Fixed from Pilot Timings}
\label{sec:m2}

\textbf{Budget clock.} The cluster's shared storage delivers at most $\sim$98 training samples per second into one
node and stalls for minutes at a time, while TFNO and AViT consume 116 and 330 per second of GPU time. A
wall-clock budget would partly measure the storage, so the 1- and 4-hour budgets count \emph{active GPU training
time}; data waits and the GPU time each job held are reported separately (Appendix~\ref{app:accounting}).

\textbf{Learning rate and batch.} An equal-budget search (20 minutes of training per run, validation only) over
$\{10^{-4}, 5\cdot10^{-4}, 10^{-3}\}$ chose $10^{-3}$ for every model (Appendix~\ref{app:more}). It is the top of the
grid, so it is the best rate tested rather than an established optimum. Batch 64, as in The Well, fits every model
in fp32 on an 80\,GB A100. TFNO's FFTs require fp32 in this implementation, so all models train in fp32; this
favours TFNO, since bf16 would speed up a training step 1.7 times for CNextU-Net and 2.2 times for AViT.

\section{Milestone 3: Results at 10\% Data}
\label{sec:results}

\textbf{One-step accuracy} (Table~\ref{tab:test}). After four hours, CNextU-Net is clearly the most accurate in
developed convection (0.050, against 0.075 for AViT and 0.109 for TFNO), while AViT is best during onset and,
narrowly, over all windows. Seeds agree to within 0.006. Persistence beats every model during onset but is 2--4.5
times worse in developed flow. The released CNextU-Net, trained on all the data, is better in developed flow
(0.038), but it also differs in schedule and hardware: by our estimate it completed at most $\sim$25 of its 200
scheduled epochs, so its learning rate never annealed. After one hour the order changes: AViT leads over all
windows (0.39) and CNextU-Net is last (0.73) and the most seed-dependent.

\begin{table}[!htbp]
\centering\small
\begin{tabular}{@{}lcccc@{}}
\toprule
& \multicolumn{3}{c}{Physical-field one-step VRMSE} & Paper protocol \\
\cmidrule(lr){2-4}
Model (10\% data, checkpoint) & all windows & onset & developed & (all windows) \\
\midrule
CNextU-Net, 4 h & 0.136 $\pm$ 0.006 & 0.387 $\pm$ 0.025 & \textbf{0.050} $\pm$ 0.000 & 0.349 $\pm$ 0.022 \\
TFNO, 4 h & 0.209 $\pm$ 0.001 & 0.502 $\pm$ 0.003 & 0.109 $\pm$ 0.000 & 0.488 $\pm$ 0.001 \\
AViT, 4 h & \textbf{0.134} $\pm$ 0.001 & \textbf{0.308} $\pm$ 0.003 & 0.075 $\pm$ 0.000 & \textbf{0.298} $\pm$ 0.002 \\
\midrule
CNextU-Net, 1 h & 0.728 $\pm$ 0.176 & 2.618 $\pm$ 0.681 & \textbf{0.081} $\pm$ 0.003 & 2.345 $\pm$ 0.600 \\
TFNO, 1 h & 0.473 $\pm$ 0.078 & 1.505 $\pm$ 0.305 & 0.120 $\pm$ 0.000 & 1.362 $\pm$ 0.249 \\
AViT, 1 h & \textbf{0.394} $\pm$ 0.036 & \textbf{1.299} $\pm$ 0.140 & 0.085 $\pm$ 0.000 & \textbf{1.171} $\pm$ 0.125 \\
\midrule
Released CNextU-Net (100\%, paper) & 0.224 & 0.766 & 0.038 & 0.669 \\
Released TFNO (100\%, paper) & 0.300 & 0.623 & 0.190 & 0.663 \\
Persistence & 0.199 & 0.121 & 0.226 & 0.202 \\
\bottomrule
\end{tabular}
\caption{One-step test error, mean $\pm$ standard deviation over two seeds (0.000: below 0.0005). Bold: best of our
models at each budget. The released TFNO lacks its Fourier layers (Appendix~\ref{app:tfno}).}
\label{tab:test}
\end{table}

\textbf{Rollouts} (Table~\ref{tab:rollout}, Fig.~\ref{fig:rollouts}). One-step accuracy does not carry over in the
same order. After four hours CNextU-Net is the most stable: its error after 30 steps from developed flow (0.53,
0.40) is below persistence's (0.86, 0.54), and 89\% of its rollouts stay below VRMSE 1. AViT, despite its one-step
error, degrades fastest; after 30 steps from developed flow it is no better than persistence (0.72 against 0.70).
After one hour TFNO is the most stable (60\% of rollouts below VRMSE 1, against 34\% for the others), and no model
beats persistence after 30 steps. Persistence keeps the most rollouts below the thresholds only because its error
levels off near 0.7; over the first steps every model is 2--5 times more accurate (Fig.~\ref{fig:rollouts}a).

\begin{table}[!htbp]
\centering\footnotesize
\begin{tabular}{@{}lccc@{}}
\toprule
& VRMSE after 30 steps & Fraction staying below threshold & Paper protocol, steps \\
Model & from frame 20 / 80 / 140 & through step 30: 1 / 2 / 5 & 6--12 / 13--30 (normalised) \\
\midrule
CNextU-Net, 4 h & 1.78 / \textbf{0.53} / \textbf{0.40} & \textbf{0.89} / \textbf{0.94} / \textbf{0.98} & \textbf{15} / \textbf{30} \\
TFNO, 4 h & 1.79 / 0.59 / 0.45 & 0.75 / 0.85 / 0.91 & 51 / 77 \\
AViT, 4 h & 5.42 / 0.80 / 0.65 & 0.63 / 0.80 / 0.87 & 75 / 167 \\
\midrule
CNextU-Net, 1 h & 3.60 / 1.09 / 1.04 & 0.34 / 0.74 / 0.79 & 202 / 148 \\
TFNO, 1 h & \textbf{2.67} / \textbf{0.90} / \textbf{0.78} & \textbf{0.60} / \textbf{0.77} / \textbf{0.86} & \textbf{81} / \textbf{116} \\
AViT, 1 h & 4.67 / 1.11 / 0.98 & 0.34 / 0.74 / 0.83 & 129 / 184 \\
\midrule
Released CNextU-Net (100\%) & & & 36 / 35 \\
Released TFNO (100\%) & & & 142 / 169 \\
Persistence & 0.68 / 0.86 / 0.54 & 0.94 / 1.00 / 1.00 & 0.27 / 0.50 \\
\bottomrule
\end{tabular}
\caption{Rollouts on the test set (physical fields unless stated; mean over two seeds; 525 rollouts per model).
Right: the paper's Table 3 protocol, rollouts from frame 0 (published: $>10$). Bold: best of our models at each
budget.}
\label{tab:rollout}
\end{table}

\begin{figure}[!htbp]
\centering
\figureorpending{fig_rollouts.pdf}{scale=1}
\caption{Four-hour models and persistence. (a) Rollout VRMSE per step from frames 80 and 140 (developed
convection). (b) Fraction of all 525 rollouts still below VRMSE 1 at each step (Kaplan--Meier survival).}
\label{fig:rollouts}
\end{figure}

\textbf{Accuracy against cost} (Fig.~\ref{fig:cost}, Table~\ref{tab:resources}). From one to four hours the
developed-flow error falls by 38\% for CNextU-Net, 12\% for AViT and 9\% for TFNO. To \emph{train}, AViT is the
most efficient on the GPU (six times as many optimizer steps as CNextU-Net, a seventh of the memory), but it reads
data faster than our storage delivers, so each four-hour run held a GPU for 11--12 hours, against 5--6 for
CNextU-Net. To \emph{use}, every model forecasts 30 steps in at most a quarter of a second with under 0.5\,GB of
memory. For one forecast at a time TFNO is the fastest (135\,ms) and CNextU-Net the slowest (249\,ms); for many
forecasts at once AViT is the cheapest per sample (0.89\,ms per step at batch 64, against 3.15 for TFNO and 6.74
for CNextU-Net).

\begin{figure}[!htbp]
\centering
\figureorpending{fig_accuracy_cost.pdf}{width=\linewidth}
\caption{Accuracy against cost, 10\% data (seed mean; bars: seed range). (a) Developed-flow one-step VRMSE at one
and four hours. (b) The same four-hour error against the GPU time each run held, including data waits. (c) Error
after 30 steps from frames 80 and 140 against the time for one 30-step forecast (batch 1). Dashed: released
CNextU-Net (100\% data); dotted: persistence.}
\label{fig:cost}
\end{figure}

\begin{table}[!htbp]
\centering\footnotesize
\setlength{\tabcolsep}{4.5pt}
\begin{tabular}{@{}lcccccccc@{}}
\toprule
& \multicolumn{4}{c}{Training (4-hour budget, batch 64)} & \multicolumn{4}{c}{Inference} \\
\cmidrule(lr){2-5}\cmidrule(l){6-9}
& Optimizer & & Peak memory, & GPU time & per step, & 30 steps, & per sample, & Memory, \\
Model & steps & Epochs & alloc.\ (reserved) & held per run & batch 1 & batch 1 & batch 64 & batch 1 \\
\midrule
CNextU-Net & 12,122--12,194 & 28 & 60.9 (64.8)\,GB & 5.2--6.0\,h & 8.3\,ms & 249\,ms & 6.74\,ms & 0.31\,GB \\
TFNO & 26,339--26,352 & 61 & 48.6 (56.9)\,GB & 6.9--7.1\,h & 4.4\,ms & 135\,ms & 3.15\,ms & 0.44\,GB \\
AViT & 74,978--75,060 & 175 & 8.8 (8.9)\,GB & 11.1--12.2\,h & 5.4\,ms & 163\,ms & 0.89\,ms & 0.23\,GB \\
\bottomrule
\end{tabular}
\caption{Cost to train and to use on one A100 80\,GB, fp32. GPU time held includes data waits, start-up,
validation and checkpoints. Per sample at batch 64: forward pass in the training profile, an upper bound for
inference.}
\label{tab:resources}
\end{table}

\textbf{Physical fidelity} (Figs.~\ref{fig:spectra} and~\ref{fig:snapshots}). Per wavenumber, the median ratio of
predicted to true kinetic energy over the 35 regimes shows where each model is wrong. CNextU-Net stays within
0.80--1.57 below $k_x=100$. TFNO is too smooth: between $k_x\approx25$ and 100 its ratio falls to 0.33 after one
step and 0.09 after 30. AViT is close after one step but gains energy over a rollout (0.55--21 below $k_x=100$
after 30 steps), and both CNextU-Net and AViT put far too much energy in the smallest scales, where the true energy
is tiny. The snapshots agree: CNextU-Net follows the plumes through 30 steps, TFNO blurs them and builds a warm
layer along the bottom wall, and AViT shows patch-sized rectangular artifacts by step 12.

\begin{figure}[!htbp]
\centering
\figureorpending{fig_spectra.pdf}{scale=1}
\caption{Horizontal kinetic-energy spectra of the four-hour models. (a) One regime ($\mathrm{Ra}=10^8$,
$\mathrm{Pr}=1$) after 30 steps. (b, c) Predicted over true energy after 1 and 30 steps: median over the 35
regimes (band: middle 50\%); 1 is exact. Shaded: small scales ($k_x>100$).}
\label{fig:spectra}
\end{figure}

\begin{figure}[!htbp]
\centering
\figureorpending{fig_snapshots_buoyancy_rayleigh_benard_Rayleigh_1e8_Prandtl_1__0__s080.pdf}{scale=1}
\caption{Buoyancy at $\mathrm{Ra}=10^8$, $\mathrm{Pr}=1$, rolled out from frame 80 by our four-hour models (seed 0)
after 1, 12 and 30 steps. One colour range for all panels. $\mathrm{Ra}=10^{10}$ in Appendix~\ref{app:more}.}
\label{fig:snapshots}
\end{figure}

\section{Why the Families Behave Differently}
\label{sec:why}

The results follow from how each architecture represents the flow, and they match known behaviour of neural PDE
surrogates.

\textbf{Spectral truncation sets TFNO's ceiling.} A Fourier layer mixes information over the whole domain, but only
through the lowest modes, here 16 in each direction ($k_x\lesssim25$); finer scales pass only through pointwise
layers~\cite{li2021fno,kossaifi2024tfno}. This bias suits smooth, large-scale dynamics, consistent with TFNO's
early lead in rollouts after one hour. It also sets a floor that more training does not lower: TFNO's training loss
stops falling after about 1.5 hours, its developed-flow error improves by only 9\% from one to four hours, and it
loses most of the energy between $k_x\approx25$ and 100, blurring the thin plumes. More modes, or a hybrid with
local convolutions, would be the usual remedies~\cite{gupta2022pdearena}.

\textbf{Locality at several resolutions suits convection, once CNextU-Net is trained.} Convection is driven by
local structures (thermal boundary layers, rising and falling plumes) organised into large rolls, which matches a
U-Net's local convolutions at four resolutions. U-Net variants are among the strongest baselines in neural PDE
benchmarks~\cite{gupta2022pdearena,ohana2024well}. CNextU-Net is also the most expensive per step (61\,GB, 12,000
optimizer steps in four hours against AViT's 75,000), so it needs the longest training to pay off: least reliable
after one hour, best after four, and still improving fastest.

\textbf{The patch sets AViT's resolution.} AViT turns $16\times16$-pixel patches into tokens and mixes them by
attention~\cite{mccabe2024mpp}. That is cheap (six times as many steps per hour) and gives the best one-step error
over all windows, but detail inside a patch is decoded patch by patch. Errors at the patch scale, visible as
rectangular artifacts and as excess small-scale energy, then grow over a rollout.

\textbf{One-step training does not guarantee stable rollouts.} All three models minimise the one-step MSE. That
loss weights each scale by its energy, so low-energy small scales are barely constrained (the spectral bias of
neural networks~\cite{rahaman2019spectral}), and surrogates trained this way are known to neglect non-dominant,
often high-frequency, components~\cite{lippe2023pderefiner}. In a rollout the model then receives its own slightly
wrong outputs, inputs it never saw in training~\cite{brandstetter2022mp}, and uncontrolled growth of such errors is
a known route to instability~\cite{mccabe2023stability}. This is why one-step rankings do not carry over to
rollouts, and why the excess small-scale energy of CNextU-Net and AViT matters. Training on the model's own outputs
(the pushforward trick~\cite{brandstetter2022mp}), injecting noise or filtering high frequencies address it; they
are outside this study.

\textbf{Persistence and the metric.} With frames 0.25 time units apart, the flow changes little per step, so
copying the last frame is a strong baseline. VRMSE is relative to the target's spatial variance, which is near zero
during onset, so small absolute errors there become large relative ones. A useful model must beat persistence where
the flow evolves; ours do so by 2--4.5 times per step in developed convection, but not during onset.

\textbf{Data or compute?} After 28 to 175 epochs over 140 trajectories, every model's validation error is still
falling at the end of training (Appendix~\ref{app:more}), so none shows signs of overfitting within four hours: at
this budget, compute rather than data appears to be the binding constraint. The 100\% runs will test this directly.

\section{Findings Under Limited Data and Compute}
\label{sec:findings}

These hold for these configurations, 10\% of the data and the tested budgets (Table~\ref{tab:choice}).
\begin{itemize}
    \item \textbf{Accuracy.} After four hours CNextU-Net is the most accurate in developed convection (0.050) and
    the most stable over rollouts (89\% stay below VRMSE 1); AViT matches it over all windows but degrades fastest;
    TFNO is too smooth.
    \item \textbf{One hour.} AViT has the lowest one-step error and TFNO the most stable rollouts; CNextU-Net is the
    least reliable, but gains most from more training (38\% against 12\% and 9\%).
    \item \textbf{Cost.} A four-hour run held a GPU for 5--6 (CNextU-Net), 7 (TFNO) or 11--12 hours (AViT,
    storage-bound). A single 30-step forecast takes 135--249\,ms, fastest with TFNO; in batches AViT is the
    cheapest per forecast.
    \item \textbf{Negative results.} No model beats persistence during onset; after 30 steps from developed flow
    AViT is no better than persistence; TFNO's error stops improving early.
    \item \textbf{Caveats.} Two seeds on one fixed subset; fp32 favours TFNO; the learning rate comes from a short
    search.
\end{itemize}

\begin{table}[!htbp]
\centering\footnotesize
\begin{tabularx}{\linewidth}{@{}>{\raggedright\arraybackslash}p{0.27\linewidth}>{\raggedright\arraybackslash}p{0.2\linewidth}>{\raggedright\arraybackslash}X@{}}
\toprule
If the priority is & Best of the three & Trade-off \\
\midrule
Accurate forecasts of developed convection that stay stable over 30 steps & CNextU-Net & Slowest single forecast
(249\,ms per 30 steps); 61\,GB training memory at batch 64; needs the longest training \\ \addlinespace[3pt]
Fastest single forecast (batch 1) & TFNO (135\,ms per 30 steps) & Too smooth (loses energy between $k_x\approx25$
and 100); stops improving early \\ \addlinespace[3pt]
Least training memory, cheapest batched forecasts & AViT (8.8\,GB; 0.89\,ms per sample at batch 64) & Needs
storage delivering $\sim$330 samples/s; rollouts degrade fastest \\ \addlinespace[3pt]
Only one hour of training & AViT (one step); TFNO (rollouts) & CNextU-Net varies most between seeds; none beats
persistence after 30 steps \\ \addlinespace[3pt]
Convection onset (first 50 frames) & None: persistence & Every model is less accurate than copying the last frame
\\
\bottomrule
\end{tabularx}
\caption{Which model to use under which constraint, from the 10\% results after four hours unless stated.}
\label{tab:choice}
\end{table}

\section{Revised Plan}
\label{sec:plan}

The core comparison is unchanged: three families, 10\% and 100\% of the training data, two seeds, checkpoints at
one and four hours. The 100\% runs complete the data axis: repeating Fig.~\ref{fig:cost} and Table~\ref{tab:choice}
with all the data will measure what ten times more data buys at each budget, and whether CNextU-Net's lead and
TFNO's ceiling persist. Changes from the abstract:
\begin{itemize}
    \item \textbf{Budgets} count active GPU training time; GPU time held, data waits and cost are reported
    separately.
    \item \textbf{Evaluation} separates onset from developed convection, reports all three failure thresholds, and
    adds the paper protocol for reference.
    \item \textbf{Resources.} The core study needs about 135 allocated GPU-hours instead of the abstract's 100,
    because of data waits (Appendix~\ref{app:accounting}). The 10\% work used 72.6 GPU-hours (2,160 SU billed).
    \item \textbf{Optional}, after the core: a Dedalus~\cite{burns2020dedalus} solver baseline, to set the models'
    forecast time and accuracy against the numerical solver's; a second draw of the 10\% subset; 24-hour runs with
    the paper's learning rates.
\end{itemize}
\textbf{Next steps.} (1) Make data loading robust to storage stalls and test node-local staging; (2) run the six
100\% runs and the data--compute comparison; (3) complete field-wise VRMSE and the Nusselt-number error. The
abstract's extensions (input-noise robustness, a third seed, 30\% data, ConvLSTM) follow only after the core.

\clearpage
\appendix
% Appendix headings read "Appendix A: ..."; \ref still gives the bare letter.
\makeatletter
\renewcommand{\@seccntformat}[1]{\ifcsname prefix@#1\endcsname\csname prefix@#1\endcsname\fi
  \csname the#1\endcsname\ifcsname prefix@#1\endcsname:\fi\quad}
\newcommand{\prefix@section}{Appendix~}
\makeatother

\section{Training Recipe}
\label{app:recipe}

Reference settings come from the paper and The Well's configuration at commit \texttt{e63dec8}. Our software: The
Well \texttt{e63dec8}, neuraloperator 0.3.0, PyTorch 2.8.0 (CUDA 12.8), Python 3.13.

\begin{table}[!htbp]
\centering\footnotesize
\begin{tabularx}{\linewidth}{@{}>{\raggedright\arraybackslash}p{2.0cm}>{\raggedright\arraybackslash}X>{\raggedright\arraybackslash}X@{}}
\toprule
& This study & The Well paper / code \\
\midrule
Activation & GELU (all three models) & GELU (same model code) \\
Optimizer & AdamW, $\beta = (0.9, 0.999)$, weight decay $10^{-4}$, no clipping & AdamW, weight decay $10^{-4}$ in the code (0.01 in the text), no clipping \\
Loss & one-step MSE, normalised units, equal field weights & same \\
Batch size, precision & 64, fp32 & 64, fp32 \\
Learning rate & $10^{-3}$ for all, from our search over $\{10^{-4}, 5\cdot10^{-4}, 10^{-3}\}$ & per model from $\{10^{-4}, 5\cdot10^{-4}, 10^{-3}, 5\cdot10^{-3}, 10^{-2}\}$: $5\cdot10^{-4}$ (CNextU-Net), $10^{-4}$ (TFNO) \\
Schedule & 500 warmup steps, cosine to 0 over the 4-hour budget & 5 warmup epochs, cosine to $0.1\cdot$lr over 200 epochs, stopped at 12\,h \\
Budget, hardware & 1\,h and 4\,h of active GPU training on one A100 & 12\,h on one H100, best validation epoch \\
Training data & 10\% of the trajectories (100\% next) & 100\% \\
\bottomrule
\end{tabularx}
\caption{Training recipe: ours against the paper's.}
\label{tab:recipe}
\end{table}

\textbf{Released CNextU-Net's schedule.} Our A100 trains CNextU-Net at 54 samples per second; with an H100 at most
three times faster, 12 hours give at most $\sim$7 million samples, about 25 epochs of 274,400 windows. With 5
warmup epochs and a cosine over 200, the learning rate at epoch 25 is still above 95\% of its peak. This is an
estimate; the model card does not give the epoch count.

\section{TFNO Forward Pass}
\label{app:tfno}

In The Well's code, commit \texttt{2d8d7eaf} (2024-06-03) called TFNO's Fourier blocks without keeping their
output, so the model computed only its lifting and projection layers, a per-pixel network; commit
\texttt{e4301602} (2025-11-17) restored the assignment~\cite{wellfix}. The paper's TFNO results and the released
checkpoint date from that period. Scored with the forward pass it was trained with, the released TFNO reproduces
the published value; with the corrected forward pass the same weights score $\approx$10.5. Our runs use the fixed
code. Upstream issue \#49 reports related inconsistencies~\cite{wellissue}.

\section{Metric Details}
\label{app:metrics}

\textbf{VRMSE.} The Well's \texttt{VRMSE} class at \texttt{e63dec8}, over the $N=512\times128$ points of a target
frame, with $s_f^2$ the unbiased spatial variance; the four fields are averaged with equal weights and windows with
equal weight. The paper protocol scores normalised fields $(y_f-\mu_f)/\sigma_f$ with The Well's training
statistics ($\sigma_f=0.249, 0.149, 0.208, 0.180$ for $b,p,u,w$). Without $\varepsilon$ the two would be identical,
since the ratio is invariant to an affine rescaling; with $\varepsilon$, normalising shrinks the floor from
$10^{-7}$ to the equivalent of $\varepsilon\sigma_f^2\approx2$--$6\times10^{-9}$. During onset the velocity is
nearly zero (exactly zero at $t=0$) and its spatial variance is comparable to $10^{-7}$, so the physical-field score
is damped there and the normalised one is not. This is the whole difference between the released CNextU-Net's
0.669 (paper protocol) and 0.224 (physical fields): on onset windows they score 2.510 and 0.766, while from frame
50 on both give 0.038.

\textbf{Failure horizon.} The first step whose field-mean VRMSE exceeds a threshold (1, 2 or 5). With no rollout
ending before step 30, the Kaplan--Meier survival at step 30 equals the fraction of rollouts staying below the
threshold.

\textbf{Kinetic-energy spectrum.} At each grid row the velocity is Fourier-transformed along the periodic $x$ axis;
$E(k_x)=\tfrac12(|\hat u|^2+|\hat w|^2)$ (one-sided) at $k_x=2\pi n/L_x$ with $L_x=4$, averaged over the
Chebyshev-spaced rows with quadrature weights, so that summing over $k_x$ gives the mean kinetic energy. Spectra are
averaged over each regime's 15 rollouts (5 trajectories $\times$ 3 starts) at steps 1, 10 and 30. The spectrum
describes how energy is spread over scales, not where structures are.

\textbf{Inference time.} Batch 1, fp32, input already on the GPU: after 10 warm-up predictions, the mean of 50
one-step predictions and one full 30-step rollout, each ended by a CUDA synchronisation. At batch 1 the GPU is far
from fully used, so latency includes per-call overhead.

\section{Resource Accounting}
\label{app:accounting}

\textbf{Timer.} Active training time is measured per optimizer step with Python's \texttt{perf\_counter}, from the
moment the data loader returns the batch to just after a CUDA synchronisation following the optimizer update. It
excludes waiting for data, validation, checkpoint writes and start-up. After a crash the clock resumes from the
last checkpoint, so lost steps are not counted.

\textbf{Validation sample.} 64 fixed windows evenly spaced over the validation split's 34,300 windows, covering all
35 regimes; 28\% are onset windows. It monitors training and does not replace the test evaluation.

\textbf{Billing.} The paid partition charges 32 SU per GPU-hour plus 1 SU per core-hour (40 SU per hour for one
A100 and 8 cores); the free partition is not billed. Of the 72.6 allocated GPU-hours so far, 54.2 were on the paid
partition (2,160 SU) and 18.3 on the free one.

\begin{table}[!htbp]
\centering\footnotesize
\begin{tabular}{@{}lccc@{}}
\toprule
Hours per run & CNextU-Net & TFNO & AViT \\
\midrule
\multicolumn{4}{@{}l}{\textit{10\% runs, measured}} \\
Queue wait before training started & 0.1--0.9 & 0.5--0.7 & 6.7--8.8 \\
GPU training (the budget) & 4.0 & 4.0 & 4.0 \\
Waiting for data & 0.8--1.1 & 2.5--2.8 & 6.1--7.3 \\
Start-up, and steps replayed after a preemption or crash & 0.13--0.65 & 0.14--0.15 & 0.68--0.69 \\
Validation (64 windows every 30 minutes) & 0.01--0.02 & 0.02--0.10 & 0.18 \\
Checkpoint writes & 0.23--0.24 & 0.08--0.15 & 0.03--0.13 \\
\textbf{Allocated time of training} & \textbf{5.2--6.0} & \textbf{6.9--7.1} & \textbf{11.1--12.2} \\
Test evaluation job (1-hour and 4-hour checkpoints) & 0.87--0.95 & 0.57 & 0.34--0.38 \\
SU per training run (paid partition) & 189--191 & 274--286 & 443--488 \\
\midrule
\multicolumn{4}{@{}l}{\textit{100\% runs, projected}} \\
Waiting for data & 1.0--2.0 & 3.0--4.5 & 9.5--12.5 \\
\textbf{Allocated time of training} & \textbf{5.5--6.5} & \textbf{7.5--9} & \textbf{14--17} \\
\bottomrule
\end{tabular}
\caption{Time per run (range over the two seeds). CNextU-Net's first attempts ran on the free partition, so its SU
covers 4.8\,h of its 5.2--6.0\,h. The 100\% projection assumes the measured storage rate ($\sim$98 samples/s) and
that the full training split ($\sim$290\,GB) no longer fits in a node's page cache.}
\label{tab:time}
\end{table}

\begin{table}[!htbp]
\centering\footnotesize
\begin{tabular}{@{}lr@{}}
\toprule
Component & Allocated GPU-hours \\
\midrule
Six 10\% runs, completed (24.0 of them active training) & 49.8 \\
Learning-rate search, profiling, pipeline check, environment checks, first replication try & 18.5 \\
Test evaluation of the six 10\% runs and persistence (completed) & 4.2 \\
Six 100\% runs (projected, 24.0 active) & $\sim$54--65 \\
Test evaluation of the 100\% runs & $\sim$4.5 \\
\midrule
Core study & $\sim$131--142 \\
\bottomrule
\end{tabular}
\caption{Allocated GPU-hours of the core study: completed parts from Slurm accounting, the rest projected. No
contingency is included; the abstract's $\sim$100 hours included $\sim$30.}
\label{tab:budget}
\end{table}

\begin{figure}[!htbp]
\centering
\figureorpending{fig_loader.pdf}{scale=1}
\caption{Training samples per second read from the shared file system by one node, against the rate at which each
model consumes them on the GPU (horizontal lines).}
\label{fig:loader}
\end{figure}

\FloatBarrier
\section{Additional Results}
\label{app:more}

\begin{table}[!htbp]
\centering\small
\begin{tabular}{@{}lccc@{}}
\toprule
Validation VRMSE after 20 min & $10^{-4}$ & $5\cdot10^{-4}$ & $10^{-3}$ \\
\midrule
CNextU-Net & 0.841 & 0.469 & \textbf{0.408} \\
TFNO & 0.603 & 0.358 & \textbf{0.302} \\
AViT & 0.646 & 0.332 & \textbf{0.257} \\
\bottomrule
\end{tabular}
\caption{Learning-rate search: 20 minutes of training per run, 10\% data, 512 validation windows.}
\label{tab:lr}
\end{table}

\begin{figure}[!htbp]
\centering
\figureorpending{fig_training.pdf}{scale=1}
\caption{Training MSE (normalised units, smoothed over 3 minutes) and validation VRMSE on the 64 fixed windows
(physical fields), against active training time, 10\% data. Solid: seed 0; dashed: seed 1.}
\label{fig:curves}
\end{figure}

\begin{figure}[!htbp]
\centering
\figureorpending{fig_test_by_window.pdf}{scale=1}
\caption{One-step test VRMSE (physical fields) by window: our four-hour models (seed mean; band: seed range), the
released CNextU-Net trained on all the data, and persistence. Shaded: onset windows.}
\label{fig:testwindow}
\end{figure}

\begin{figure}[!htbp]
\centering
\figureorpending{fig_snapshots_buoyancy_rayleigh_benard_Rayleigh_1e10_Prandtl_1__0__s080.pdf}{scale=1}
\caption{As Fig.~\ref{fig:snapshots}, at $\mathrm{Ra}=10^{10}$: the true flow is far finer than any model
reproduces; CNextU-Net keeps the large-scale plume pattern but smooths the small structures.}
\label{fig:snapshots10}
\end{figure}

\clearpage

\begin{thebibliography}{99}
\setlength{\itemsep}{6pt}

\bibitem{ohana2024well}
Ohana, R., et al.
``The Well: a Large-Scale Collection of Diverse Physics Simulations for Machine Learning.''
NeurIPS, Datasets and Benchmarks Track, 2024.\\
\url{https://arxiv.org/abs/2412.00568}

\bibitem{liu2022convnext}
Liu, Z., Mao, H., Wu, C.-Y., Feichtenhofer, C., Darrell, T., and Xie, S.
``A ConvNet for the 2020s.''
CVPR, 2022.\\
\url{https://arxiv.org/abs/2201.03545}

\bibitem{li2021fno}
Li, Z., Kovachki, N., Azizzadenesheli, K., Liu, B., Bhattacharya, K., Stuart, A., and Anandkumar, A.
``Fourier Neural Operator for Parametric Partial Differential Equations.''
ICLR, 2021.\\
\url{https://arxiv.org/abs/2010.08895}

\bibitem{kossaifi2024tfno}
Kossaifi, J., Kovachki, N., Azizzadenesheli, K., and Anandkumar, A.
``Multi-Grid Tensorized Fourier Neural Operator for High-Resolution PDEs.''
TMLR, 2024.\\
\url{https://arxiv.org/abs/2310.00120}

\bibitem{mccabe2024mpp}
McCabe, M., et al.
``Multiple Physics Pretraining for Spatiotemporal Surrogate Models.''
NeurIPS, 2024.\\
\url{https://arxiv.org/abs/2310.02994}

\bibitem{gupta2022pdearena}
Gupta, J. K., and Brandstetter, J.
``Towards Multi-spatiotemporal-scale Generalized PDE Modeling.''
arXiv:2209.15616, 2022.\\
\url{https://arxiv.org/abs/2209.15616}

\bibitem{rahaman2019spectral}
Rahaman, N., Baratin, A., Arpit, D., Draxler, F., Lin, M., Hamprecht, F. A., Bengio, Y., and Courville, A.
``On the Spectral Bias of Neural Networks.''
ICML, 2019.\\
\url{https://arxiv.org/abs/1806.08734}

\bibitem{lippe2023pderefiner}
Lippe, P., Veeling, B. S., Perdikaris, P., Turner, R. E., and Brandstetter, J.
``PDE-Refiner: Achieving Accurate Long Rollouts with Neural PDE Solvers.''
NeurIPS, 2023.\\
\url{https://arxiv.org/abs/2308.05732}

\bibitem{brandstetter2022mp}
Brandstetter, J., Worrall, D., and Welling, M.
``Message Passing Neural PDE Solvers.''
ICLR, 2022.\\
\url{https://arxiv.org/abs/2202.03376}

\bibitem{mccabe2023stability}
McCabe, M., Harrington, P., Subramanian, S., and Brown, J.
``Towards Stability of Autoregressive Neural Operators.''
TMLR, 2023.\\
\url{https://arxiv.org/abs/2306.10619}

\bibitem{burns2020dedalus}
Burns, K. J., Vasil, G. M., Oishi, J. S., Lecoanet, D., and Brown, B. P.
``Dedalus: A flexible framework for numerical simulations with spectral methods.''
Physical Review Research 2, 023068, 2020.\\
\url{https://arxiv.org/abs/1905.10388}

\bibitem{wellrb}
The Well, \texttt{rayleigh\_benard} dataset card, commit \texttt{e63dec8}.\\
\url{https://github.com/PolymathicAI/the_well/blob/e63dec8/datasets/rayleigh_benard/README.md}

\bibitem{wellfix}
The Well source code, commits \texttt{2d8d7eaf9eb46b96f215abf39c94b6be4bf9ea5c} (2024-06-03) and
\texttt{e4301602e57f2dddeaa87856dda49bbba47c663b} (2025-11-17).\\
\url{https://github.com/PolymathicAI/the_well/commit/2d8d7eaf9eb46b96f215abf39c94b6be4bf9ea5c}\\
\url{https://github.com/PolymathicAI/the_well/commit/e4301602e57f2dddeaa87856dda49bbba47c663b}

\bibitem{wellissue}
The Well, issue \#49, ``Inconsistent Pre-trained Baseline Performance'' (open since 2025-04-21).\\
\url{https://github.com/PolymathicAI/the_well/issues/49}

\end{thebibliography}

\end{document}
